% arXiv requires \pdfoutput=1 within the first few lines when using pdflatex.
\pdfoutput=1

\documentclass[11pt]{article}

% ---------------------------------------------------------------------------
% Page and typography
% ---------------------------------------------------------------------------
\usepackage[margin=1in]{geometry}
\usepackage[T1]{fontenc}
\usepackage[utf8]{inputenc}
\usepackage{microtype}
\usepackage{lmodern}

% ---------------------------------------------------------------------------
% Math, tables, figures
% ---------------------------------------------------------------------------
\usepackage{amsmath}
\usepackage{amssymb}
\usepackage{amsthm}
\usepackage{booktabs}
\usepackage{graphicx}
\usepackage{subcaption}
\usepackage{xcolor}
\usepackage{array}

% ---------------------------------------------------------------------------
% Algorithms
% ---------------------------------------------------------------------------
\usepackage{algorithm}
\usepackage{algpseudocode}

% ---------------------------------------------------------------------------
% Bibliography. natbib + bibtex is the safest choice for arXiv, which does not
% run bibtex: you must upload the generated .bbl alongside the sources.
% ---------------------------------------------------------------------------
\usepackage[numbers,sort&compress]{natbib}
\bibliographystyle{plainnat}

% ---------------------------------------------------------------------------
% Links and cross-references (hyperref last, cleveref after it)
% ---------------------------------------------------------------------------
\usepackage[hidelinks,breaklinks]{hyperref}
\usepackage[capitalise,noabbrev]{cleveref}

% ---------------------------------------------------------------------------
% Code listings, with a YAML dialect since listings has no built-in one
% ---------------------------------------------------------------------------
\usepackage{listings}

\definecolor{lstkey}{HTML}{0B5394}
\definecolor{lstcomment}{HTML}{6B7280}
\definecolor{lststring}{HTML}{7A3E9D}
\definecolor{lstbg}{HTML}{F7F7F8}

\lstdefinelanguage{yaml}{
  keywords={true,false,null,apiVersion,kind,metadata,spec,status},
  keywordstyle=\color{lstkey}\bfseries,
  sensitive=true,
  comment=[l]{\#},
  commentstyle=\color{lstcomment}\itshape,
  morestring=[b]", % chktex 18 (listings config, not prose)
  stringstyle=\color{lststring},
}

\lstset{
  basicstyle=\ttfamily\footnotesize,
  backgroundcolor=\color{lstbg},
  breaklines=true,
  captionpos=b,
  columns=fullflexible,
  frame=none,
  keepspaces=true,
  showstringspaces=false,
  xleftmargin=1em,
  aboveskip=0.8em,
  belowskip=0.8em,
}

\lstnewenvironment{yamllisting}[1][]
  {\lstset{language=yaml,#1}}
  {}

% ---------------------------------------------------------------------------
% Project macros. Use these everywhere so terminology stays consistent and a
% late rename is a one-line change.
% ---------------------------------------------------------------------------
\newcommand{\sysname}{\textsc{agent-tasks}\xspace}
\usepackage{xspace}

\newcommand{\crd}[1]{\texttt{#1}}
\newcommand{\AgentTask}{\crd{AgentTask}}
\newcommand{\AgentStep}{\crd{AgentStep}}
\newcommand{\AgentWorkflow}{\crd{AgentWorkflow}}
\newcommand{\AgentStepClass}{\crd{AgentStepClass}}
\newcommand{\AgentClusterPolicy}{\crd{AgentClusterPolicy}}
\newcommand{\TrackerConnection}{\crd{TrackerConnection}}
\newcommand{\Sandbox}{\crd{Sandbox}}

% The pure convergence core is two functions (implementation.md, section 11).
\newcommand{\fold}{\textsc{fold}\xspace}
\newcommand{\decide}{\textsc{decide}\xspace}
% Math symbols: no \xspace, so write \snapshot{} where prose follows.
\newcommand{\snapshot}{\ensuremath{h}}
\newcommand{\gatekey}{\ensuremath{\kappa}}
\newcommand{\ledger}{\ensuremath{\mathcal{L}}}

% Outline convention used in the section stubs:
%   [impl]  shipped by the Phase 1 implementation (implementation.md)
%   [spec]  specified there but delivered in a later phase
% The prose must not present a [spec] item as if it were running.

% Drafting aids. Comment out the definitions below before submission.
\newcommand{\todo}[1]{\textcolor{red}{\textbf{[TODO: #1]}}} % chktex 13
\newcommand{\note}[1]{\textcolor{blue}{\textbf{[NOTE: #1]}}} % chktex 13
% \renewcommand{\todo}[1]{}
% \renewcommand{\note}[1]{}

\theoremstyle{definition}
\newtheorem{definition}{Definition}
\theoremstyle{plain}
\newtheorem{proposition}{Proposition}

% ---------------------------------------------------------------------------
% Front matter
% ---------------------------------------------------------------------------
\title{\sysname: Content-Addressed Verification for Autonomous\\
       Coding Agents on Kubernetes}

% TODO: confirm author list, affiliations, and ORCIDs before submission.
\author{
  Author One\thanks{Primitive. Correspondence: \texttt{author@example.com}} \\
  \and
  Author Two \\
}

\date{\today}

\begin{document}

\maketitle

\begin{abstract}
  % Write this LAST. 150-250 words. Structure:
%   1. Autonomous coding agents produce code faster than humans can review it.
%   2. Existing orchestration assumes an immutable tree; agents mutate it mid-run.
%   3. We present a content-addressed verification model + a Kubernetes operator.
%   4. Key mechanisms: (gate, tree-hash) ledger, path-scoped invalidation,
%      capability/composition split, microVM isolation, post-merge validation loop.
%   5. Result: one measured datapoint on AWS. Fill from Section 13.
%   6. Open source, Apache-2.0.

\todo{Write abstract last, after Sections 6, 9, and 13 are stable.}

\end{abstract}

\section{Introduction}
\label{sec:intro}

% 1.5-2 pages.

\subsection{Problem}
% The unit of work is a ticket; the desired output is a merged pull request.
% Producing the diff is increasingly tractable. Establishing that the diff is
% correct is not, and it is what actually gates deployment.

\subsection{Why existing orchestration does not fit}
% Two families, both mismatched:
%   - CI/CD and workflow engines (Argo, Tekton, Airflow, Temporal): assume the
%     tree is immutable for the duration of the run. Nodes execute once.
%   - Agent frameworks (LangGraph, OpenHands, SWE-agent): model the reasoning
%     loop but provide no isolation, no capacity model, no cluster primitives.
% Neither treats verification state as a first-class, re-entrant artifact, and
% neither separates the step that changes the tree from the step that judges it.

\subsection{Contributions}
% Enumerate 5. Draft:
%   (C1) A formulation of agent scaling as a verification-capacity problem, with
%        verification density as the tunable parameter.
%   (C2) A content-addressed verification ledger with default-closed,
%        path-scoped invalidation and a pure convergence core (\fold, \decide),
%        giving well-defined convergence and a computable termination bound for
%        a tree that mutates between gate runs.
%   (C3) A declarative workflow specification that separates capability
%        (cluster admin: what a step IS -- action or gate) from composition
%        (team: which steps, in what order, invalidated by what), with
%        integrity gates injected by policy at plan resolution so that no
%        workflow can compose them away or shadow them with a weaker copy.
%   (C4) A threat model for agent-authored code in a shared cluster, including
%        gate-gaming as adversarial selection pressure under best-of-N and the
%        observation that invalidation defaults are a security property.
%   (C5) An open-source Kubernetes implementation -- five CRDs, two controllers,
%        no admission webhooks -- and a measured cost datapoint.

\todo{Draft after Section 2 and Section 6 settle.}

\section{The Specification-Only Organization}
\label{sec:thesis}

% ~1.5 pages. This is the load-bearing framing section. Write it early.

% Thesis sentence (settled, refine wording):
%   Autonomous agent systems scale not by removing human judgment but by
%   relocating it from implementation review, which is O(agent output), to
%   specification and post-merge intent validation, which is O(tickets).
%   Verification density determines how much of the gap between those two the
%   system can close autonomously.

\subsection{Verification capacity as the bound}
% At N=1 a human reviews every diff and review IS the trust mechanism.
% At N=1000 that is arithmetically impossible. Trust must come from executable
% verification. Gates are not CI hygiene; they are the substitute for attention.

\subsection{Two checks, neither substitutable}
% Gates verify correctness against specification.
% Humans verify specification against intent.
% Only the first is automatable. This division is the contribution's real shape.

\subsection{Where the work goes}
% Table: Product / QA / Engineer, before and after. Nobody writes application
% code; everybody writes specification and verification. This is a relocation of
% expertise, not a reduction -- state that explicitly to avoid overclaiming.

\subsection{Verification compounds; review does not}
% A human review does not accumulate. A gate, once written, applies to every
% subsequent task forever. Findings become tests (see Section 12). This is the
% mechanism by which trustworthiness compounds and N scales.

\subsection{The autonomy dial}
% Two knobs, per workflow, per risk class:
%   escalationPolicy: Fail | Retry | Elicit    (default Fail -- asking does not scale)
%   mergePolicy:      Manual | AutoOnGreen | AutoWithSampledAudit
% Position on the dial is a function of measured verification density, not a
% philosophical stance. Interactive mode is a workflow-development affordance.

\subsection{Trust is purchased with compute}
% Raising verification density multiplies compute across tasks x gates x
% candidates. Verification is the CHEAP axis (elastic, shared-kernel) which is
% why the substitution is economically viable at all. Forward-reference 8.6.

\section{Motivation}
\label{sec:motivation}

% 2-3 pages.

\subsection{Anatomy of the convergence loop}
% Walk a real ticket: edit, lint, typecheck, unit, integration, review.
% 50-200 edit-run cycles. Emphasise: the tree mutates between every gate run.
% Name the two kinds of step the loop contains and keep them apart from here
% on: ACTIONS change the tree (the agent's edit, a formatter's fix); GATES judge
% it (lint, typecheck, tests). An action is done or not done; a gate's verdict
% is about one snapshot and no other.

\subsection{Why one-shot execution fails}
% The core observation. A DAG engine node fires once when dependencies are
% satisfied. A gate must fire repeatedly as the tree changes.
% NOTE: do NOT frame this as "graphs vs DAGs" -- our dependency graph IS acyclic
% and admission-enforced to be so (CEL rules plus controller-side plan
% resolution; there are no admission webhooks). The contrast is EXECUTION
% SEMANTICS:
%   one-shot topological execution  vs.  content-addressed level-triggered convergence.
% Table: topology / trigger / node identity / completion criterion.

\subsection{Why agent frameworks fail}
% No kernel isolation, no capacity model, no multi-tenancy, no declarative
% workflow surface, no tracker-neutral integration.

\subsection{Requirements}
% R1 verification results must be re-entrant and content-keyed
% R2 invalidation must be path-scoped or cost grows superlinearly -- and it
%    must fail closed, or scoping becomes a way to skip verification
% R3 agent-authored artifacts must not run on shared kernels
% R4 the agent must not be able to weaken its own verification (this becomes:
%    policy-owned gates, class-only behaviour, spec the agent cannot write)
% R5 capacity is finite and non-elastic at the isolation layer
% R6 tracker integration must not leak into the core
% R7 the decision procedure must be pure: testable without a cluster and
%    replayable after a controller crash

\section{Related Work}
\label{sec:related}

% 1.5-2 pages. MANDATORY for arXiv credibility.

\subsection{Workflow and pipeline engines}
% Argo Workflows, Tekton, Airflow, Temporal, Make.
% Argo is the closest comparison: has memoization, but keyed on a user-supplied
% cache key, not content-addressed with cascading invalidation. Acyclic by
% definition; loops require recursion templates. Node status hits the etcd
% object limit, motivating status offloading -- evidence for our design: we
% keep only the bounded working set the reconciler needs in status and project
% history out (Section 8.4), so the limit is designed around rather than hit.

\subsection{Build systems: our actual lineage}
% Bazel, Nix, Buck, Nx. Build Systems a la Carte for the formal vocabulary.
% Claim the lineage explicitly: our ledger is a build system's incrementality
% model applied to a tree that an agent mutates rather than a human. The gate
% key includes the resolved verifier, so redefining a gate invalidates its
% results exactly as changing a compiler invalidates its objects.
% Difference: build systems assume the tree is fixed during a build.

\subsection{Agent frameworks and benchmarks}
% LangGraph, OpenHands, SWE-agent, SWE-bench.
% These model the reasoning loop; we model the infrastructure and verification
% substrate underneath it. Complementary, not competing.

\subsection{Isolation}
% gVisor, Kata Containers, Firecracker, agent-sandbox, commercial sandboxes.
% We consume agent-sandbox rather than reimplementing; position \sysname as the
% task-orchestration layer directly above it.

\subsection{Kubernetes extension patterns}
% Operator pattern, CRDs, CEL validation rules as the webhook-free admission
% path, Kueue for admission, level-triggered reconciliation.

\section{Design Goals and Non-Goals}
\label{sec:goals}

% 0.5-1 page. Short, sharp, referenced back to throughout.

\subsection{Goals}
% G1 Convergence over execution: terminate at a verified fixed point.
% G2 Declarative and Kubernetes-native: the CRDs are the contract.
% G3 Capability/composition separation enforced by RBAC and admission.
% G4 Tracker-neutral core: zero tracker-specific code in the operator.
% G5 Standalone operation: kubectl apply with no external dependency -- and no
%    admission webhooks; the install is CRDs plus one manager.
% G6 Sublinear verification cost via path-scoped invalidation, without ever
%    trading a false carry-forward for it.
% G7 Determinism and replayability: a task is admitted once against a
%    complete, immutable resolved plan; every decision is a pure function of
%    recorded state and survives a controller crash unchanged.

\subsection{Non-goals}
% N1 We do not model the agent's internal reasoning loop.
% N2 We do not provide the isolation runtime (delegated to agent-sandbox).
% N3 We are not a general workflow engine; acyclic one-shot DAGs are better
%    served by existing tools, and we use one as a step executor.
% N4 Determining WHAT to build is the adopting organisation's product process
%    and is out of scope. We assume a ticket exists and is specifiable.

\section{The Verification Ledger}
\label{sec:ledger}

% 3-4 pages. THE core section. Write this first, or second after Section 2.
% Normative source: implementation.md section 2 (semantics) and section 11
% (executable contract). Every claim here must trace to one of those.

\subsection{Definitions}
% Snapshot \snapshot: immutable content identity of the workspace (git tree
%   hash; submodule and LFS pointers are tree content).
% Action: a step that may mutate the workspace. Completing it yields a new
%   snapshot or leaves the snapshot unchanged. An action has run-state, never
%   validity -- it is not "green forever".
% Gate g: a deterministic predicate over a snapshot, producing pass/fail plus
%   evidence. A gate never mutates the workspace.
% Driver: the one action the workflow designates as the agent. Runs once at
%   task start, and again whenever a failed gate needs a fix that no cheaper
%   fix action resolved.
% Fix action: an ordinary action class referenced by a gate class; tried once
%   per snapshot before a driver run is spent.
% Gate key \gatekey = hash(snapshot, resolved gate specification): changing
%   the verifier invalidates its old results, exactly as in a build system.
% Ledger \ledger: partial function (g, \gatekey) -> {pass, fail(e)}.
%   Kubernetes materialises it as a BOUNDED view (current result and verified
%   snapshot per gate, explicit counters), not an append-only log (Section 8.4).
% Formalise: one \begin{definition} block per item; Appendix B carries the
% formal statements, keep this subsection readable.
% The sentence to land: actions mutate state; gates assert properties of state.
% Everything in the model follows from keeping those two apart.

\subsection{Path-scoped invalidation}
% Each gate declares invalidatedBy: a set of path globs.
% On an edit producing h' from h, gate g retains its result iff
%   diff(h, h') INTERSECT invalidatedBy(g) = empty set.
% THE DEFAULT IS CLOSED: an undeclared invalidatedBy means ["**"] -- any change
% invalidates. Carry-forward is an explicit optimisation the author opts into.
% Argue the asymmetry: a false carry-forward is a correctness and security bug;
% an unnecessary re-run is merely expensive. (Section 9.3 shows how the open
% default would have neutralised the injected integrity gate.)
% Failed results carry forward under the same rule as passed ones; a truncated
% or contradictory change report is treated as ["**"].
% Two gate scopes: tree-scoped gates key on the snapshot alone; diff-scoped
% gates (the suppression meta-gate, Section 9) key on (base revision, snapshot)
% [spec].
% Corollary worth stating: a rebase is just another tree change. It invalidates
% diff-scoped gates and preserves tree-scoped ones whose paths it misses. This
% is what makes high monorepo concurrency affordable -- and it is why the base
% revision must NOT be folded into snapshot identity itself.

\subsection{\fold and \decide as pure functions}
% Two functions, no Kubernetes dependency, no I/O, no clock:
%   \fold(view, execution result) -> view: adopts the snapshot, applies
%     invalidation, updates counters, detects oscillation.
%   \decide(plan, view, budgets) -> RunAction | RunGate | Escalate | Done.
% Include the pseudocode (implementation.md section 11.3). Points to make:
%   - cost-ordered, fail-fast: steps walk in dependency order; a cheap failing
%     gate short-circuits before an expensive one runs;
%   - fix action before driver: remediation escalates in cost;
%   - budgets guard REMEDIATION, never verification -- the last budgeted fix is
%     always observed before any escalation;
%   - blocking is honoured: a non-blocking failure is recorded and never blocks
%     the fixed point;
%   - the driver is a step: a task whose base tree already passes every gate
%     still runs it. Verifying the base tree is never success.
% Unit-testable against a synthetic ledger. The implementation's test matrix
% (implementation.md section 11.7) is the evidence: cite case counts and the
% property suite, not adjectives.

\subsection{Termination}
% Fixed point: every blocking gate green against a SINGLE snapshot AND the
%   driver completed at least once.
% Oscillation: a newly produced snapshot equals a non-parent ancestor in the
%   lineage => the loop is cycling; break immediately. A no-op action (the
%   snapshot is unchanged) is NOT oscillation.
% A lifetime cap on driver runs plus a per-gate cap on fix attempts replace
%   the sliding-window budget of earlier drafts: with one driver the window
%   bought nothing but complexity.
% Global budgets: tokens; ACTIVE wall clock (accumulates from provisioning,
%   pauses while suspended); consecutive infrastructure failures, escalating
%   only where no fallback exists (driver, blocking gates) -- an exhausted fix
%   action degrades to the driver, an exhausted non-blocking gate is tolerated.
% Diff-size ratchet (monotonic growth with gates red => flailing, revert): an
%   early-exit heuristic layered on termination, not part of it [spec].
% \begin{proposition}: bounded termination with the explicit bound of
%   implementation.md section 11.5 -- driver and fix runs bound snapshots,
%   snapshots bound gate runs, consecutive-reset bounds errors. Proof sketch in
%   Appendix B.

\subsection{Comparison to one-shot execution}
% Table from Section 3.2, expanded. This is where the reader should conclude
% that a DAG engine cannot express this without becoming a ledger.
% Add the row a pipeline cannot have: "gate result becomes unknown because the
% tree changed underneath it" (line 17 of the reference trace).

\section{Declarative Workflow Specification}
\label{sec:workflow}

% 1.5-2 pages.

\subsection{Capability versus composition}
% \AgentStepClass (cluster-scoped, admin): images, commands, executor, timeout,
%   evidence format, autofix, catalog dependencies. This is CAPABILITY.
% \AgentWorkflow (namespaced, team): which classes, dependency edges,
%   invalidation globs, candidates, budgets, selection. This is COMPOSITION ONLY.
% Security consequence: a tenant authoring a workflow cannot introduce an image.
% RBAC falls out: write on workflows, read-only on classes.

\subsection{Declaring dependencies and invalidation}
% after: cost ordering and fail-fast.
% invalidatedBy: the ledger's invalidation rules, as data rather than code.
% Include the minimal and the full example listings side by side.

\subsection{Resolution and version pinning}
% Resolution order: explicit ref, tag match, repo match, sole workflow, default.
% SECURITY: never resolve from ticket free text. Body and comments are writable
% by anyone with issue access, including the agent. Resolve from trusted
% structured fields only, once, at admission, before any agent code runs.
% Pin specHash: editing a workflow mid-flight would silently void ledger entries.

\subsection{Non-composable integrity requirements}
% \AgentWorkflowDefaults injects required integrity steps via mutating webhook.
% A tenant must not be able to compose away the suppression meta-gate.
% Forward-reference Section 9.

\section{Architecture}
\label{sec:arch}

% 3-4 pages. Normative source: implementation.md sections 3, 5 and 6.

\subsection{Resource model}
% Table: kind / scope / writer / purpose.
% \AgentTask (ns; human or connector), \AgentStep (ns; operator),
% \AgentWorkflow (ns; team), \AgentStepClass (cluster; admin),
% \AgentClusterPolicy (cluster; admin -- a singleton named default [impl],
% namespace selectors arriving with multi-tenancy [spec]), \TrackerConnection
% (ns; connector phase [spec]).
% \AgentStep is a self-contained, immutable execution record: everything the
% executor needs is copied from the resolved plan into its spec, so executing
% it reads no mutable object. Its content-keyed name is derived from stored
% counters, and the name IS the idempotency key. TTL after completion is
% operator logic -- a CRD gets no ttlSecondsAfterFinished for free.

\subsection{Controllers and the convergence core}
% Two controllers. The task controller is a thin shell around two pure
% functions: it feeds completed \AgentStep results to \fold, persists the
% bounded view, calls \decide, and materialises the answer. Single-flight
% [impl]: one execution at a time, so a mutation never lands under an
% in-flight gate; parallel gates are a later phase [spec].
% CRITICAL: the operator must not run the agent inside Reconcile. An LLM turn is
% a long, expensive, non-idempotent side effect; reconcile is level-triggered.
% Executor contract: Ensure(executionID, spec) / Poll / Cancel. Ensure is
% idempotent by construction -- the controller may call it any number of times
% -- which makes crash recovery natural rather than bolted on. Contrast with
% dispatch-then-record-a-handle, which is at-least-once however carefully the
% handle is guarded.
% Lifecycle: success cleanup is ordinary reconciliation (Validated -> release
% -> Succeeded, after the work namespace is observed gone); the finalizer is
% the deletion path only. Suspend cancels the in-flight execution and retains
% workspace and view; the wall-clock budget pauses.

\subsection{Isolation and placement}
% One work namespace per task: the isolation and networking boundary for
% everything derived from the agent's workspace (default-deny network
% policies, quota, permitted runtime classes, materialised from the policy
% template). Created at provisioning, released on success, retained for
% post-mortem on failure, torn down explicitly on deletion -- a namespaced
% object cannot own a Namespace, so none of this is garbage-collected for free.
% microVM (Kata/Firecracker) on bare metal for anything derived from the
% agent's workspace; pinned-digest catalog images on elastic shared-kernel
% nodes. Gates run against read-only workspace mounts, with the executor
% echoing the input snapshot as a backstop.
% PROVENANCE RULE, enforced by ValidatingAdmissionPolicy:
%   images matching sandbox/* MUST set runtimeClassName kata-* and live in a
%   work namespace.
% Note the filesystem constraint that drives co-location: inotify does not cross
% NFS, and Firecracker is block-only, so the agent and its dev servers share one
% microVM rather than a distributed filesystem.

\subsection{State: hot path and follower}
% etcd holds the control plane and the ledger's BOUNDED working set: current
% snapshot, the full-but-bounded snapshot lineage, per-gate result and
% verified snapshot, explicit counters (never derived from history), spend,
% the one active execution. It is the ONLY thing on the reconcile path.
% Nothing authoritative in two places: status is authoritative for
% reconciliation; a separate projector appends facts to Postgres for history,
% evidence and the read API [spec]; connectors render from status to trackers.
% Postgres down => history stops accruing, pipeline keeps running. Not tier-1.
% Contrast with Argo, whose offloaded status IS on the reconcile path.
% Gap-free projection: AgentStep finalizer released only after the append commits.

\subsection{Capacity and admission}
% Two scarce resources with different shapes: sandbox slots (per task, hours,
% metal-bound) and inference (per step, seconds, rate-limit-bound).
% Slot tenure is the dominant lever: commit-per-step makes the sandbox
% disposable, enabling release on wait, preemption at step boundaries, and crash
% recovery through one mechanism.
% Capacity gating is separate from the controller workqueue.

\subsection{Cost considerations}
\label{sec:cost}
% ~1 page. Qualitative framing + ONE measured datapoint. Explicitly hedged.
% Three axes and how they scale:
%   C_sandbox    ~ concurrent_tasks x tenure           (poor elasticity, metal)
%   C_verify     ~ tasks x affected_gates x candidates (good elasticity)
%   C_inference  ~ turns                               (rented variable / owned fixed)
% Which mechanism attacks which term:
%   tenure          <- slot release, commit-per-step
%   affected_gates  <- path-scoped invalidation, affected-target scoping
%   candidates      <- pressure-sensitive N
%   inference       <- self-hosted GPU converts variable to fixed
% The closed invalidation default spends verification compute for safety;
% report its measured overhead against an open default as the price of the
% invariant, not as waste.
% State plainly: absolute figures depend on instance mix, repo size, gate suite,
% model choice, and rented-vs-owned inference. Numbers are a single-configuration
% datapoint, not a general model. Readers should reproduce on their workload.
% Then one table (see Section 13 for the measurement) and one interpretation
% paragraph. Write the interpretation LAST.
\todo{Fill the cost table after the AWS experiment run. Do not draft the
  interpretation paragraph against placeholder numbers.}

\section{Threat Model}
\label{sec:threat}

% 2-3 pages. Most novel material in the paper. Consider cross-listing cs.CR.

\subsection{Trust boundaries}
% What kernel isolation buys: protection of the host kernel and the rest of the
% cluster from an escaping agent.
% What it does NOT buy: protection from anything deliberately granted. In
% practice, damage looks like correct use of granted authority -- force-push
% over main, dropping a table believed to be a fixture, rewriting a CI workflow,
% exfiltration through an allowed egress path.

\subsection{Gate-gaming}
% The agent will make gates pass by disabling them. Not maliciously: it is the
% locally optimal move. Enumerate: noqa, type-ignore, eslint-disable, skip
% markers, deleted assertions, ignore-file additions, lowered thresholds,
% --no-verify.
% The suppression meta-gate: runs on the DIFF, not the tree. Rejects gate-config
% modification, net-new suppressions, test deletion, assertion-count decrease,
% gate binary/version changes.

\subsection{Best-of-N as adversarial selection pressure}
% KEY INSIGHT. With N=1 you get an agent that occasionally games a gate. With
% N=5 you explicitly SELECT the candidate that gamed most effectively, because
% gaming is the cheapest path to green.
% Therefore the meta-gate must run BEFORE ranking, and trip => disqualify, not
% penalise. Note that a "smaller diff" tiebreak makes this worse in isolation:
% a type-ignore comment is a very small diff.

\subsection{Confused deputy}
% Co-located tool pods launder egress: CDP on a browser pod is arbitrary code
% execution against a less-isolated peer; a credential-holding git pod is worse.
% CI is the sharpest edge: an agent that can modify workflow definitions runs
% code on trusted infrastructure with trusted credentials, and the microVM is
% irrelevant at that point.
% IMDS: cloud metadata endpoints must be unreachable at both the policy layer
% and the instance-metadata-options layer.

\subsection{Capability scoping}
% Tools expose intent-level APIs, not raw protocol access.
% The credential never enters the sandbox: the agent commits locally and the
% broker PULLS, so a compromised sandbox has nothing to steal.
% Repo-scoped, short-lived tokens with minimal permissions.

\subsection{Residual risks}
% Be honest. Post-merge validation (Section 12) is the external check that
% gaming cannot reach, but it is sampled, not exhaustive.

\section{Multi-Candidate Execution}
\label{sec:candidates}

% 1-1.5 pages. Merge into Section 6 if page budget is tight.

\subsection{Diversity and failure independence}
% Resampling at higher temperature is weak: same solution, different names.
% Stronger: different strategy hints, different models, different context scope.
% Fan-out only buys variance reduction. If candidates fail identically because
% the mistake is implied by the prompt or the repo's existing patterns, N=3
% costs 3x and catches nothing. This is measurable -- report the correlation.

\subsection{Selection}
% Gates filter; they do not rank. Tiebreak order:
%   1. meta-gate clean (disqualifying, not penalising)
%   2. no new dependencies
%   3. smaller diff
%   4. reviewer score

\subsection{Why hunk-level composition fails}
% "Take the best parts of each" fails on SEMANTIC conflict, not textual: three
% diffs may all apply cleanly and produce something no agent understands.
% It also voids the ledger -- the composed tree was validated by nobody.
% What works instead: pick a winner, feed the losers' diffs as review input,
% one agent integrates, gates re-run on a normal tree hash.

\subsection{Fan out cheap, converge before expensive}
% Fan out across phase-0 gates (in-sandbox, no catalog dependencies).
% Converge to one candidate BEFORE provisioning stateful dependencies.
% Reviewers and research agents need no sandbox at all -- they do not execute
% code, so fanning them out costs tokens, not scarce isolation capacity.

\section{Issue Tracker Integration}
\label{sec:tracker}

% 1-1.5 pages.

\subsection{A provider-neutral status vocabulary}
% status.plan, status.question, status.conditions render as a Linear agent plan,
% a Jira checklist comment, or a GitHub task list.
% Design rule: if a status field only makes sense because one tracker has a
% corresponding concept, it belongs in the connector, not the core.

\subsection{One-directional at both edges}
% webhook -> connector -> AgentTask.spec
% AgentTask.status -> connector (watch) -> tracker
% The operator never calls the tracker. A connector crash is harmless: on
% restart it re-projects from current status. No replay logic.

\subsection{Webhook plus reconciling sweep}
% Webhooks are lossy. The sweep is not hygiene -- it is what stops abandoned
% work from holding scarce capacity. Drift cases: missed delegation, ticket
% closed mid-run, delegation removed, failed projection.

\subsection{Interactive mode}
% Elicitation as a workflow-development affordance, not a core state. Default
% escalation policy is Fail, because a system that scales by asking more
% questions re-centralises the bottleneck it exists to remove.

\section{The Post-Merge Validation Loop}
\label{sec:validation}

% ~1.5 pages. This section carries the thesis empirically. Do not cut it.
% STATUS: [spec] -- lands with merge policies and the tracker connectors.

\subsection{Intent validation is above the gaming layer}
% Gates verify correctness against specification; they are attackable by the
% thing being gated. A human confirming "this does what the ticket asked" is not
% reading the diff -- they are testing behaviour against the requirement.
% Suppression comments and weakened assertions do not survive that check.
% Crucially it is O(tickets), not O(agent-turns): a task that took 200 turns and
% 3 candidates gets exactly one validation.

\subsection{Failure is a normal outcome}
% Post-merge validation failure must be as cheap and structured as a gate
% failure: revert, new task carrying the finding, and -- the important part --
% the finding becomes a gate so it cannot recur. Concretely: a new gate class,
% made mandatory through the cluster policy's required gates, injected into
% every plan admitted from then on (Section 7.5).

\subsection{Compounding trust}
% This loop is the mechanism by which the system becomes more trustworthy over
% time. It is also what justifies raising a repo's position on the autonomy dial.

\subsection{Sampled audit}
% Audit k\% of auto-merged PRs; track the rate at which audits catch what gates
% missed. Constant cost, not linear in task count. This is the empirical trust
% measure and the basis for the mergePolicy decision.

\section{Preliminary Evaluation}
\label{sec:eval}

% 2-3 pages. Without numbers this paper gets skimmed. Even preliminary helps.

\subsection{Setup}
% Full disclosure or the section is decorative: instance types, region,
% on-demand vs spot, model and version, gate suite, N, run window (AWS and
% model pricing both move). One repo -- self-hosting on the \sysname repo is a
% legitimate and appealing framing. 20-50 tickets across at least two cost
% classes so the spread is visible.
% REPORT THE FAILURES: escalated tasks, post-merge rejections, meta-gate
% rejections. Cost per merged PR that silently excludes failed attempts is
% misleading and reviewers check for exactly this.

\subsection{Convergence behaviour}
% Turns per gate. Escalation reason distribution (oscillation-heavy => gates
% conflict; budget-heavy => ladder too slow). Gate evidence quality as a
% confound: bad error output is often the actual bug.

\subsection{Invalidation efficiency}
% Post-rebase gate survival rate. Gate-runs avoided by path scoping, as a
% percentage. These validate the central claim of Section 6.

\subsection{Cost}
% Decomposed by axis, not a single dollar figure. Feeds the table in 8.6.
% Transferable ratios (usable by readers with different pricing):
%   slot utilisation = active step seconds / slot-held seconds
%   verification:generation compute ratio
%   tokens per merged PR

\subsection{Safety}
% Meta-gate rejection rate. Whether it rises with N -- the prediction from 9.3.
% Audit-catch rate.

\subsection{Trust over time}
% THE THESIS AS A CURVE:
%   post-merge validation failure rate should decline as gates accumulate
%   human interventions per merged PR should decline while merge rate holds
%   gate-suite growth vs escape rate: does converting findings to gates actually
%     reduce recurrence, or do failures keep arriving in novel forms? If the
%     latter, verification density has a ceiling and we must say so.

\subsection{Threats to validity}
% Single repo, single organisation, small N, short window, one model family.

\section{Limitations and Future Work}
\label{sec:limitations}

% ~1 page.

\subsection{Limitations}
% - Verification density has an unknown ceiling; some correctness properties
%   resist expression as gates.
% - Sampled audit is not exhaustive; gaming that survives both gates and
%   sampling is undetected by construction.
% - The closed invalidation default trades compute for safety. Narrow declared
%   scopes recover most of it -- but only where authors declare them.
% - The implementation is single-candidate and single-flight: one execution at
%   a time, no parallel gates. Both are specified, neither is built.
% - Bare-metal capacity is the hard bound and provisions slowly.
% - Single-controller throughput; sharding is future work.
% - Evaluation is one repo, one organisation, one model family.

\subsection{Scope boundary}
% Determining WHAT to build is the adopting organisation's product process. We
% assume a ticket exists and is specifiable. State this as an assumption, not a
% weakness -- it is where the system's contract begins.

\subsection{Future work}
% - Parallel gate execution: an instance dimension in the ledger key, and one
%   driver run addressing every failure at a snapshot rather than one per cycle.
% - Multi-candidate execution (Section 10) with per-candidate ledgers.
% - Automatic gate synthesis from post-merge findings.
% - Self-hosted inference: the crossover point where owned GPU capacity beats
%   rented tokens is determined by utilisation and is computable from traces.
% - Cross-cluster federation; sharding by namespace.
% - Formal treatment of convergence under adversarial gate-gaming.

\section{Conclusion}
\label{sec:conclusion}

% 0.5 page. Restate the thesis, the mechanism, the result. No new material.
\todo{Write last.}

\section*{Availability}
\addcontentsline{toc}{section}{Availability}

An open-source implementation of \sysname is available at
\url{https://github.com/primitivecorp/agent-tasks} under the Apache-2.0
license. Reference manifests reproducing the configurations described in
\cref{sec:workflow} are included in \cref{app:manifests}.

% Deployment instructions belong in the repository README, not in the paper.
% A how-to section reads as documentation and weakens the submission.


\bibliography{refs}

\appendix
\section{Reference Manifests}
\label{app:manifests}

% Complete, runnable YAML for:
%   - AgentStepClass: a gate with a fixAction, the fix action it names, the
%     driver (an Action class), and the integrity gate the policy requires
%   - AgentWorkflow (minimal and full)
%   - AgentClusterPolicy (the singleton named default)
%   - AgentTask (standalone, no tracker source)
%   - RuntimeClass with scheduling and overhead
%   - The provenance ValidatingAdmissionPolicy
% Keep these in sync with the repo's test/e2e/testdata so they stay runnable.
% Show the resolved plan for the full workflow once, so the reader sees the
% injected <name>@policy gate and the synthesised <gate>@fix action that never
% appear in the authored YAML.

\todo{Populate from repository testdata once Phase 1 is green.}

\section{Invalidation Semantics}
\label{app:semantics}

% Formal statements for Section 6, following implementation.md section 11 (the
% executable contract). Definitions first, then the fixed point, then the
% termination proposition whose proof Section 6.4 defers here.

\begin{definition}[Snapshot]
  A snapshot \snapshot{} is the content identity of the workspace: the git
  tree hash, with submodule and LFS pointers taken as tree content. Snapshots
  are immutable. An action either produces a new snapshot or leaves the
  current one unchanged.
\end{definition}

\begin{definition}[Gate key]
  For a gate \(g\) with resolved specification \(s_g\) and a snapshot
  \snapshot{}, the gate key is \(\gatekey(g, \snapshot) = H(\snapshot, s_g)\)
  for a collision-resistant hash \(H\). A diff-scoped gate additionally
  includes the base revision \(b\): \(\gatekey(g, \snapshot) = H(b, \snapshot,
  s_g)\).
\end{definition}

\begin{definition}[Ledger]
  The ledger is a partial function
  \(\ledger : (g, \gatekey) \rightharpoonup \{\textsc{pass}, \textsc{fail}(e)\}\),
  where \(e\) is evidence. Absence is \textsc{unknown}, and \textsc{unknown}
  is what schedules a gate.
\end{definition}

\begin{definition}[Invalidation]
  Let an action transform \snapshot{} into \(\snapshot'\) with changed path
  set \(P\), and let \(I(g)\) be the glob set declared by gate \(g\), which
  defaults to \(\{\texttt{**}\}\). The result of \(g\) carries forward to
  \(\snapshot'\) if and only if \(P \cap \mathrm{paths}(I(g)) = \emptyset\);
  otherwise it is absent at \(\snapshot'\). If \(P\) is reported truncated,
  \(P\) is taken to be every path. The rule applies to \textsc{fail} results
  exactly as to \textsc{pass}.
\end{definition}

\begin{definition}[Fixed point]
  Let a task have driver \(d\), blocking gates \(B\) and non-blocking gates
  \(N\). The task is at its fixed point at snapshot \snapshot{} if and only if
  \(d\) has completed at least once, every \(g \in B\) has
  \(\ledger(g, \gatekey(g, \snapshot)) = \textsc{pass}\), and every
  \(g \in N\) either has a result at \snapshot{} or has exhausted its
  infrastructure retries.
\end{definition}

\begin{definition}[Oscillation]
  Let \(\snapshot_0, \ldots, \snapshot_k\) be the lineage of distinct
  snapshots a task has occupied, with \(\snapshot_0\) the base. An action
  producing \(\snapshot' \neq \snapshot_k\) oscillates if and only if
  \(\snapshot' = \snapshot_i\) for some \(i < k\). An action producing
  \(\snapshot' = \snapshot_k\) is a no-op and is not oscillation.
\end{definition}

\begin{proposition}[Bounded termination]
  Let \(R\) be the cap on driver runs, \(F\) the number of gates that name a
  fix action, \(A\) the per-gate cap on fix attempts, \(G\) the number of
  gates, and \(E\) the cap on consecutive infrastructure failures of one step.
  Every task reaches its fixed point or escalates after at most
  \[
    \bigl( (R + FA) + G\,(R + FA + 1) \bigr)\,(E + 1) + c
  \]
  executions, for a small constant \(c\).
\end{proposition}

\begin{proof}[Proof sketch]
  Each decision either terminates (\textsc{done}, \textsc{escalate}) or
  schedules one execution that strictly consumes a bounded quantity. Driver
  completions are bounded by \(R\) and fix completions by \(FA\), so completed
  actions number at most \(R + FA\). Each new snapshot requires a completed
  action, so the lineage has at most \(R + FA + 1\) entries. A gate re-runs
  only after a new snapshot has cleared its result, so gate completions are
  bounded by \(G\,(R + FA + 1)\). An infrastructure failure increments a
  per-step counter that only a success of that step resets, so each step
  suffers at most \(E\) failures per success and \(E\) more before escalating,
  which gives the factor \(E + 1\). Oscillation, the wall-clock budget and the
  token budget only tighten the bound. The implementation's property suite
  asserts this arithmetic directly against generated execution sequences.
\end{proof}

\todo{Tighten the constant \(c\) once the property tests pin it.}


\end{document}
